% =============================================================
% paper/preamble.tex —— 论文全局宏包、格式、符号宏、自定义命令
% 唯一真源：docs/NOTATION.md（符号）、paper/refs.bib（参考文献）
% 编译方式：XeLaTeX（latexmk -xelatex），严禁 pdflatex（中文会炸）
% =============================================================

% ---------- 页面设置（数模论文常规 2.5cm 页边距） ----------
\usepackage[a4paper,margin=2.3cm,top=2.4cm,bottom=2.4cm]{geometry}

% ---------- 数学 ----------
\usepackage{amsmath}
\usepackage{amssymb}
\usepackage{bm}   % 向量/矩阵加粗用 \bm

% ---------- 表格（一律 booktabs 三线表，禁用竖线） ----------
\usepackage{booktabs}
% 论文含十余张表与十余幅图，统一收紧浮动体间距以控制篇幅
\setlength{\floatsep}{4pt plus 1pt minus 1pt}
\setlength{\textfloatsep}{6pt plus 1pt minus 1pt}
\setlength{\intextsep}{5pt plus 1pt minus 1pt}
\setlength{\abovecaptionskip}{3pt}
\setlength{\belowcaptionskip}{0pt}
% 放宽浮动体页面配额：本文含 15 幅图与 17 张表，若限制过严会在页面底部留下大片空白
\renewcommand{\topfraction}{0.92}
\renewcommand{\bottomfraction}{0.75}
\renewcommand{\textfraction}{0.08}
\renewcommand{\floatpagefraction}{0.80}
\setcounter{topnumber}{3}
\setcounter{bottomnumber}{2}
\setcounter{totalnumber}{4}
\usepackage{longtable}  % 跨页符号表
\usepackage{multirow}
% 假设条目：不依赖 enumitem 的可选参数（曾被 enit@postlabel 的括号陷阱绊住），
% 直接用零间距自定义环境，保证任何 TeX 发行版下都能编译。
\newenvironment{assumptionlist}{%
  \begin{list}{}{\setlength{\leftmargin}{2.4em}\setlength{\labelwidth}{2.0em}%
    \setlength{\labelsep}{0.4em}\setlength{\itemindent}{0pt}%
    \setlength{\itemsep}{1pt}\setlength{\parsep}{0pt}\setlength{\topsep}{2pt}}%
}{\end{list}}
\usepackage{array}
\usepackage{graphicx}   % \resizebox 用于超宽表格

% ---------- 颜色（xcolor 需早于 listings 加载） ----------
\usepackage{xcolor}

% ---------- 插图 ----------
\graphicspath{{figs/}}

% ---------- 流程图与示意图（TikZ；不依赖外部图片文件） ----------
\usepackage{tikz}
\usetikzlibrary{arrows.meta,positioning,shapes.geometric,fit,backgrounds,calc}
\tikzset{
  box/.style={draw, rounded corners=2pt, align=center, font=\small,
              minimum height=7mm, inner sep=2.4pt, fill=blue!4},
  data/.style={box, fill=green!8},
  phys/.style={box, fill=orange!10},
  num/.style={box, fill=blue!8},
  deliv/.style={box, fill=purple!8},
  ver/.style={box, fill=red!7},
  flow/.style={-{Latex[length=2mm]}, thick},
  note/.style={font=\scriptsize, align=left, text=black!70},
}

% ---------- 代码（Python 样式；统一 UTF-8，XeLaTeX 下中文正常） ----------
\usepackage{listings}
\lstset{
  language=Python,
  basicstyle=\ttfamily\scriptsize,
  extendedchars=true,
  inputencoding=utf8,
  showstringspaces=false,
  columns=fullflexible,
  keepspaces=true,
  frame=single,
  framexleftmargin=2pt,
  breaklines=true,
  breakatwhitespace=false,
  postbreak=\mbox{\textcolor{gray}{$\hookrightarrow$}\space},
  commentstyle=\color{gray},
  keywordstyle=\color{blue!60!black},
  stringstyle=\color{orange!80!black},
  numbers=left,
  numberstyle=\tiny\color{gray},
  tabsize=4,
}

% ---------- 数学单位（siunitx v3 语法：\qty{}{} 与 \unit{}，勿用 v2 的 \SI{}{}） ----------
\usepackage{siunitx}
\sisetup{per-mode=symbol, detect-all}

% ---------- 算法（algorithm2e；备选 algorithm+algpseudocode，二选一勿同开） ----------
\usepackage[linesnumbered,ruled]{algorithm2e}

% ---------- 参考文献（GB/T 7714 数字顺序制，中国数模论文通用格式） ----------
% main.tex 用 \bibliographystyle{gbt7714-numerical} + \bibliography{refs}。
\usepackage[sort&compress]{natbib}
\bibliographystyle{gbt7714-numerical}
\setcitestyle{super,comma}   % 上标数字引用，正文用 \cite{...}
\renewcommand{\bibfont}{\small}

% ---------- 超链接与交叉引用（hyperref 必须最后加载，cleveref 在其之后） ----------
\usepackage[hidelinks,bookmarksnumbered=true]{hyperref}
\hypersetup{
  pdftitle={圆柱形中药材热风烘干过程的传热-传质耦合模型},
  pdfauthor={},
}
\usepackage[nameinlink]{cleveref}

% =============================================================
% 符号宏（唯一真源：docs/NOTATION.md §0「三层命名映射」）
% 规则：论文里不要直接写 $T$，要写 \T；改名只改这一处，全文自动同步。
% 说明：\t 覆盖 LaTeX 原生的 tie 重音（\renewcommand）；
%       r/k/h/L/D 等常用简写已按 NOTATION.md 规避：
%       r→\rpos、k→\kcond、h→\hconv、L→\Llen、D→\Ddiff。
% =============================================================
\renewcommand{\t}{\ensuremath{t}}           % 时间
\newcommand{\T}{\ensuremath{T}}             % 药材温度
\newcommand{\C}{\ensuremath{C}}             % 水分浓度（干基含水率）
\newcommand{\rpos}{\ensuremath{r}}          % 到药材中心的径向距离
\newcommand{\R}{\ensuremath{R}}             % 药材半径（问题 4 中随时间变化）
\newcommand{\Rzero}{\ensuremath{R_0}}       % 初始半径
\newcommand{\Llen}{\ensuremath{L}}          % 药材长度
\newcommand{\cp}{\ensuremath{c_p}}          % 定压比热容
\newcommand{\kcond}{\ensuremath{k}}         % 导热系数
\newcommand{\Ddiff}{\ensuremath{D}}         % 水分扩散系数
\newcommand{\hconv}{\ensuremath{h}}         % 对流换热系数
\newcommand{\hmass}{\ensuremath{h_m}}       % 对流传质系数
\newcommand{\Tinf}{\ensuremath{T_\infty}}   % 烘房（环境）温度
\newcommand{\Cinf}{\ensuremath{C_\infty}}   % 烘房（环境）水分浓度
\newcommand{\wfrac}{\ensuremath{w}}         % 水分质量分数 w = C/(1+C)
\newcommand{\dt}{\ensuremath{\Delta t}}     % 时间步长
\newcommand{\dr}{\ensuremath{\Delta r}}     % 空间步长
\newcommand{\Nr}{\ensuremath{N_r}}          % 径向网格节点数
\newcommand{\tend}{\ensuremath{t_{\mathrm{end}}}}  % 烘干结束时刻
% （\rho、\xi 为 LaTeX 原生，无需定义）

% =============================================================
% 自定义命令
% =============================================================
% \instruction{...}：醒目的 TODO 占位提示（待填内容用）
\newcommand{\instruction}[1]{%
  \par\medskip\noindent
  \fcolorbox{orange!70!black}{orange!10}{%
    \parbox{\dimexpr\linewidth-2\fboxsep-2\fboxrule}{%
      \textbf{\textcolor{orange!60!black}{【TODO 占位】}}\, #1}}%
  \par\medskip
}

% \code{...}：行内代码
\newcommand{\code}[1]{\texttt{#1}}

% \geninput{tab/foo}{补充说明}：安全引入由 scripts 生成的片段。
% 依据 AGENTS.md 铁律 10「不许静默降级」：片段缺失时**不静默跳过**，
% 而是打印醒目的占位框，让缺表在 PDF 里一眼可见。
% \detokenize 把路径中的下划线还原为普通字符 —— 否则"文件缺失"提示本身
% 会因为 \texttt{tab/foo_bar} 里的下划线进入数学模式而报错（实测踩过）。
\newcommand{\geninput}[2]{%
  \IfFileExists{#1.tex}{%
    \input{#1}%
  }{%
    \instruction{【生成产物缺失】未找到 \texttt{\detokenize{#1}.tex}。#2
      请先运行 \texttt{make tables}（表格与数值宏）或 \texttt{make figs}（插图）。}%
  }%
}

% \genfig[宽度]{figs/xxx.pdf}{题注}{label}：安全插图（缺文件时打印占位框）
\newcommand{\genfig}[4][0.48\linewidth]{%
  \begin{figure}[!ht]
    \centering
    \IfFileExists{#2}{%
      \includegraphics[width=#1]{#2}%
    }{%
      \instruction{【插图缺失】未找到 \texttt{\detokenize{#2}}。请先运行 \texttt{make figs}。}%
    }%
    \caption{#3}\label{#4}%
  \end{figure}%
}

% \genfigwide{...}：整幅宽度（宽图专用）
\newcommand{\genfigwide}[3]{%
  \begin{figure}[!ht]
    \centering
    \IfFileExists{#1}{%
      \includegraphics[width=\linewidth]{#1}%
    }{%
      \instruction{【插图缺失】未找到 \texttt{#1}。请先运行 \texttt{make figs}。}%
    }%
    \caption{#2}\label{#3}%
  \end{figure}%
}

% 数值宏：由 make tables 从 results/ 生成，缺失时显示 ?? 而不是编译中断
\newcommand{\genvalues}[1]{%
  \IfFileExists{tab/#1.tex}{\input{tab/#1}}{%
    \PackageWarning{paper}{tab/#1.tex 缺失，对应数值宏显示为 ??}%
  }%
}
\newcommand{\qonevalues}{\genvalues{q1_values}}
\newcommand{\qallvalues}{\genvalues{problem_values}}

% 关键词（若 ctexart 已定义则不覆盖）
\providecommand{\keywords}[1]{%
  \par\medskip\noindent\textbf{关键词：}#1\par}

% =============================================================
% 版式：摘要独占一页；正文页码从摘要页起为 1（format2026.doc 第三条）
% =============================================================
\setlength{\parskip}{1pt}
\linespread{0.98}
\ctexset{
  section/format = \raggedright\bfseries\zihao{4},
  subsection/format = \raggedright\bfseries\zihao{-4},
  subsubsection/format = \raggedright\bfseries\zihao{-4},
}
% 参考文献标题用「参考文献」而不是「References」
\renewcommand{\refname}{参考文献}
